\documentclass{article}
\usepackage{amsmath,amssymb,booktabs,geometry,tabularx}
\geometry{margin=1in}



% no indent
\setlength\parindent{0pt}





\title{Normalization from 1NF to 5NF:\\
Summary Table and Step--by--Step Examples}
\author{}
\date{}















\begin{document}
\maketitle

\section*{Summary Table of Normal Forms}

\begin{table}[h]
  \centering
  \small
  \begin{tabularx}{\textwidth}{@{}l l X X X@{}}
    \toprule
    \textbf{NF} & \textbf{What It Removes}
                & \textbf{Bad Pattern (Easy Notation)}
                & \textbf{Rule (When OK)}
                & \textbf{Fix (Decomposition Pattern)} \\
    \midrule
    1NF &
    Non--atomic or multi--valued attributes &
    (No FD issue; problem is value form) &
    All attributes must be atomic. &
    Split repeating groups:
    $R \rightarrow R_1, R_2$. \\
    \addlinespace[0.3em]
    2NF &
    Partial dependency on a composite key &
    $(A + B) \to C$ but actually $A \to C$,
    where $C$ is non--key. &
    No non--key attribute depends on a proper subset
    of a candidate key. &
    For $R(A,B,C)$:
    $R \rightarrow R_1(A,C),\ R_2(A,B)$. \\
    \addlinespace[0.3em]
    3NF &
    Transitive dependency on a key &
    $A \to B,\ B \to C$ with $C$ non--key. &
    For every FD $X \to Y$: $X$ is a key
    or each attribute in $Y$ is prime. &
    For $R(A,B,C)$:
    $R \rightarrow R_1(A,B),\ R_2(B,C)$. \\
    \addlinespace[0.3em]
    BCNF &
    Any FD whose left side is not a key &
    $X \to Y$ but $X$ is not a key. &
    Every determinant must be a superkey. &
    $R \rightarrow R_1(X,Y),\ R_2(\text{rest})$. \\
    \addlinespace[0.3em]
    4NF &
    Independent multivalued dependencies &
    $A \twoheadrightarrow B$ and
    $A \twoheadrightarrow C$ with $B$ independent of $C$. &
    For every nontrivial MVD $A \twoheadrightarrow B$,
    $A$ must be a superkey. &
    $R(A,B,C) \rightarrow R_1(A,B),\ R_2(A,C)$. \\
    \addlinespace[0.3em]
    5NF &
    Join dependencies not implied by keys &
    Nontrivial join dependency
    $\text{JD}^\ast(R_1,R_2,R_3)$ with no key in any part. &
    Every nontrivial join dependency must include
    a projection containing a key. &
    $R \rightarrow R_1, R_2, R_3$. \\
    \bottomrule
  \end{tabularx}
  \caption{Normal forms 1NF--5NF: problem, bad pattern, rule, and typical fix.}
\end{table}

\clearpage
\section*{Worked Example: From Non--1NF to BCNF}

We use a single example relation that is gradually transformed.

\subsection*{Step 0: Non--1NF Table (Repeating Group)}

Consider a table \texttt{ENROLL\_RAW}:

\begin{center}
\begin{tabular}{llll}
\toprule
StudentID & StudentName & DeptName & CoursesList \\
\midrule
1 & Alice & CS   & \{CS101, CS102\} \\
2 & Bob   & MATH & \{MATH155\} \\
\bottomrule
\end{tabular}
\end{center}

Here \texttt{CoursesList} is a set of course codes in one cell, so the table
is not in 1NF.

\bigskip

\subsection*{Step 1: Convert to 1NF}

We expand the repeating group into one row per student--course pair and
also add more attributes, so our 1NF table is:

\begin{center}
\texttt{ENROLL\_1NF(StudentID, StudentName, CourseID, CourseName, DeptID, DeptName, DeptPhone)}
\end{center}

\begin{center}
\begin{tabular}{lllllll}
\toprule
StudentID & StudentName & CourseID & CourseName & DeptID & DeptName & DeptPhone \\
\midrule
1 & Alice & CS101  & Databases    & CS   & Comp.~Sci. & 555--1000 \\
1 & Alice & CS102  & Algorithms   & CS   & Comp.~Sci. & 555--1000 \\
2 & Bob   & MATH155 & Calculus I  & MATH & Mathematics & 555--2000 \\
\bottomrule
\end{tabular}
\end{center}

All attributes are atomic, so the relation is in 1NF. Assume the \textbf{primary key} is $(\text{StudentID},\text{CourseID})$.

The following functional dependencies (FDs) hold:
\begin{align*}
\text{StudentID} &\to \text{StudentName}, \\
\text{CourseID} &\to \text{CourseName}, \text{DeptID}, \\
\text{DeptID} &\to \text{DeptName}, \text{DeptPhone}, \\
(\text{StudentID},\text{CourseID}) &\to \text{all attributes.}
\end{align*}

\bigskip

\subsection*{Step 2: Remove Partial Dependencies (2NF)}

In \texttt{ENROLL\_1NF}, some non--key attributes depend on only part of
the composite key:
\begin{align*}
\text{StudentID} &\to \text{StudentName}, \\
\text{CourseID} &\to \text{CourseName}, \text{DeptID}.
\end{align*}
Both right--hand sides are non--key attributes, so the table is not in 2NF.

We remove partial dependencies by separating information about students
and courses:

\begin{itemize}
  \item \textbf{STUDENT(StudentID, StudentName)}
  \item \textbf{COURSE\_TEMP(CourseID, CourseName, DeptID, DeptName, DeptPhone)}
  \item \textbf{ENROLL\_TEMP(StudentID, CourseID)}
\end{itemize}

Sample data:

\begin{center}
\begin{tabular}{ll}
\toprule
\multicolumn{2}{c}{\textbf{STUDENT}} \\
\midrule
StudentID & StudentName \\
\midrule
1 & Alice \\
2 & Bob \\
\bottomrule
\end{tabular}
\hspace{2em}
\begin{tabular}{lllll}
\toprule
\multicolumn{5}{c}{\textbf{COURSE\_TEMP}} \\
\midrule
CourseID & CourseName & DeptID & DeptName & DeptPhone \\
\midrule
CS101  & Databases   & CS   & Comp.~Sci.   & 555--1000 \\
CS102  & Algorithms  & CS   & Comp.~Sci.   & 555--1000 \\
MATH155 & Calculus I & MATH & Mathematics  & 555--2000 \\
\bottomrule
\end{tabular}
\end{center}

\begin{center}
\begin{tabular}{ll}
\toprule
\multicolumn{2}{c}{\textbf{ENROLL\_TEMP}} \\
\midrule
StudentID & CourseID \\
\midrule
1 & CS101 \\
1 & CS102 \\
2 & MATH155 \\
\bottomrule
\end{tabular}
\end{center}

Now every non--key attribute in each relation depends on the whole key,
so these relations are in 2NF.

\bigskip

\subsection*{Step 3: Remove Transitive Dependencies (3NF)}

\textbf{COURSE\_TEMP} still has a transitive dependency:
\[
\text{CourseID} \to \text{DeptID} \to
(\text{DeptName}, \text{DeptPhone}).
\]

To reach 3NF, we split department information into its own table.

\begin{itemize}
  \item \textbf{COURSE(CourseID, CourseName, DeptID)}
  \item \textbf{DEPARTMENT(DeptID, DeptName, DeptPhone)}
\end{itemize}

Sample data:

\begin{center}
\begin{tabular}{lll}
\toprule
\multicolumn{3}{c}{\textbf{COURSE}} \\
\midrule
CourseID & CourseName & DeptID \\
\midrule
CS101  & Databases   & CS \\
CS102  & Algorithms  & CS \\
MATH155 & Calculus I & MATH \\
\bottomrule
\end{tabular}
\hspace{2em}
\begin{tabular}{lll}
\toprule
\multicolumn{3}{c}{\textbf{DEPARTMENT}} \\
\midrule
DeptID & DeptName   & DeptPhone \\
\midrule
CS   & Comp.~Sci.  & 555--1000 \\
MATH & Mathematics & 555--2000 \\
\bottomrule
\end{tabular}
\end{center}

Our full set of tables is now:
\begin{itemize}
  \item STUDENT(StudentID, StudentName)
  \item COURSE(CourseID, CourseName, DeptID)
  \item DEPARTMENT(DeptID, DeptName, DeptPhone)
  \item ENROLL\_TEMP(StudentID, CourseID)
\end{itemize}

Each FD has a key on the left--hand side, or the right--hand side is
part of a key, so the schema is in 3NF.

\subsection*{Step 4: Check BCNF}

We verify for each relation that every determinant is a superkey:
\begin{itemize}
  \item In \textbf{STUDENT}, StudentID is the key and determinant.
  \item In \textbf{COURSE}, CourseID is the key and determinant.
  \item In \textbf{DEPARTMENT}, DeptID is the key and determinant.
  \item In \textbf{ENROLL\_TEMP}, the key is
        $(\text{StudentID},\text{CourseID})$
        and there are no nontrivial FDs with a smaller determinant.
\end{itemize}

Therefore all relations are in BCNF. Our example has now been gradually
converted from a non--1NF table to BCNF.

\clearpage
\section*{4NF Example: Multivalued Dependencies}

4NF deals with independent multivalued attributes.

\subsection*{Initial 1NF Table}

Consider the relation:

\[
\text{MOVIE(Title, Actor, Language)}
\]

Sample data:

\begin{center}
\begin{tabular}{lll}
\toprule
Title & Actor & Language \\
\midrule
M1 & A1 & English \\
M1 & A1 & French \\
M1 & A2 & English \\
M1 & A2 & French \\
\bottomrule
\end{tabular}
\end{center}

For each movie:
\begin{itemize}
  \item It has a set of actors.
  \item It has a set of languages.
  \item The choice of actor is independent of the choice of language.
\end{itemize}

We have multivalued dependencies:
\[
\text{Title} \twoheadrightarrow \text{Actor},\quad
\text{Title} \twoheadrightarrow \text{Language}.
\]

The key of MOVIE is $(\text{Title}, \text{Actor}, \text{Language})$,
so \texttt{Title} is \emph{not} a superkey, and the relation violates 4NF.

\subsection*{4NF Decomposition}

We decompose along the MVDs:

\begin{itemize}
  \item \textbf{MOVIE\_ACTOR(Title, Actor)}
  \item \textbf{MOVIE\_LANG(Title, Language)}
\end{itemize}

Sample data:

\begin{center}
\begin{tabular}{ll}
\toprule
\multicolumn{2}{c}{\textbf{MOVIE\_ACTOR}} \\
\midrule
Title & Actor \\
\midrule
M1 & A1 \\
M1 & A2 \\
\bottomrule
\end{tabular}
\hspace{2em}
\begin{tabular}{ll}
\toprule
\multicolumn{2}{c}{\textbf{MOVIE\_LANG}} \\
\midrule
Title & Language \\
\midrule
M1 & English \\
M1 & French \\
\bottomrule
\end{tabular}
\end{center}

Now the only MVDs have determinants that are keys in their respective
relations, so both relations are in 4NF.

\clearpage
\section*{5NF Example: Join Dependencies}

5NF addresses join dependencies not implied by keys.

\subsection*{Initial Relation}

Consider:

\[
\text{SUPPLY(Supplier, Part, Project)}
\]

Interpretation:
Supplier $S$ can supply Part $P$ to Project $J$.

Assume the following business rules:
\begin{itemize}
  \item For each Supplier--Part pair $(S,P)$, we know whether the supplier can supply that part.
  \item For each Supplier--Project pair $(S,J)$, we know whether the supplier works on that project.
  \item For each Part--Project pair $(P,J)$, we know whether the part is used in that project.
\end{itemize}

Furthermore, we assume there is \emph{no additional constraint} on specific triples
$(S,P,J)$: whenever all three pairwise relationships hold, the triple
is valid.

The key of SUPPLY is $(\text{Supplier},\text{Part},\text{Project})$.

There is a nontrivial join dependency:
\[
\text{JD}^\ast\bigl(
  \{\text{Supplier},\text{Part}\},
  \{\text{Supplier},\text{Project}\},
  \{\text{Part},\text{Project}\}
\bigr)
\]
meaning SUPPLY can be reconstructed as the natural join of its three
binary projections.

However, none of these projections contains a key of the original
relation, so this violates 5NF.

\subsection*{5NF Decomposition}

We decompose SUPPLY into the three projections:

\begin{itemize}
  \item \textbf{SP(Supplier, Part)}
  \item \textbf{SJ(Supplier, Project)}
  \item \textbf{PJ(Part, Project)}
\end{itemize}

Sample data (illustrative):

\begin{center}
\begin{tabular}{ll}
\toprule
\multicolumn{2}{c}{\textbf{SP}} \\
\midrule
Supplier & Part \\
\midrule
S1 & P1 \\
S1 & P2 \\
S2 & P2 \\
\bottomrule
\end{tabular}
\hspace{1.5em}
\begin{tabular}{ll}
\toprule
\multicolumn{2}{c}{\textbf{SJ}} \\
\midrule
Supplier & Project \\
\midrule
S1 & J1 \\
S2 & J1 \\
\bottomrule
\end{tabular}
\hspace{1.5em}
\begin{tabular}{ll}
\toprule
\multicolumn{2}{c}{\textbf{PJ}} \\
\midrule
Part & Project \\
\midrule
P1 & J1 \\
P2 & J1 \\
\bottomrule
\end{tabular}
\end{center}

The original triples (Supplier, Part, Project) can be obtained as:
\[
\text{SUPPLY} = SP \;\Join\; SJ \;\Join\; PJ.
\]

There are no remaining nontrivial join dependencies that do not include
a projection containing a key, so this decomposition is in 5NF.

\end{document}
